\chapter{Phase 3 \& 4: Deep Residual MLP, Sign Studio and 3D Augmentations}
\label{Chapter 4}

\section{Deep Residual MLP Topology (\texttt{ISLLetterClassifier})}
To achieve high-precision classification on 126-dimensional geometric coordinate inputs with sub-2ms latency, we designed a custom \textbf{Deep Residual Multi-Layer Perceptron (\texttt{ISLLetterClassifier})}.

The network architecture comprises three dense processing blocks and a linear classification head:
\begin{enumerate}
    \item \textbf{Input Projection Block:} Projects the 126-dimensional input vector into a 256-dimensional latent space:
    \begin{equation}
        \mathbf{h}_1 = \text{Dropout}_{0.20}\left(\text{SiLU}\left(\text{BatchNorm1d}\left(\mathbf{W}_1 \mathbf{X} + \mathbf{b}_1\right)\right)\right)
    \end{equation}
    where $\mathbf{W}_1 \in \mathbb{R}^{256 \times 126}$, and $\text{SiLU}(x) = x \cdot \sigma(x)$ provides smooth non-linear gradient propagation.
    \item \textbf{Residual Bottleneck Block:} A 256-dimensional hidden layer equipped with a persistent identity skip connection:
    \begin{equation}
        \mathbf{h}_2 = \text{Dropout}_{0.20}\left(\text{SiLU}\left(\text{BatchNorm1d}\left(\mathbf{W}_2 \mathbf{h}_1 + \mathbf{b}_2\right)\right)\right) + \mathbf{h}_1
    \end{equation}
    The additive shortcut connection ($\mathbf{h}_1$) allows error gradients to backpropagate directly through the block, completely preventing gradient vanishing during multi-epoch training.
    \item \textbf{Compression Block:} Compresses the representation from 256 to 128 dimensions:
    \begin{equation}
        \mathbf{h}_3 = \text{Dropout}_{0.20}\left(\text{SiLU}\left(\text{BatchNorm1d}\left(\mathbf{W}_3 \mathbf{h}_2 + \mathbf{b}_3\right)\right)\right)
    \end{equation}
    \item \textbf{Classification Head:} A final linear projection $\mathbf{z} = \mathbf{W}_4 \mathbf{h}_3 + \mathbf{b}_4$ outputting 35 unnormalized class logits.
\end{enumerate}

\section{Loss Formulation with Label Smoothing}
To prevent overconfident predictions along ambiguous class boundaries (e.g., subtle distinctions between `M', `N', and `S' hand shapes), the network was trained using **Cross-Entropy Loss with Label Smoothing ($\alpha = 0.05$)**:
\begin{equation}
\mathcal{L}_{\text{LS}}(y, \mathbf{p}) = -\sum_{k=1}^{K} q_k \log p_k, \quad \text{where } q_k = (1 - \alpha) \cdot \mathbb{I}(y = k) + \frac{\alpha}{K}
\end{equation}
where $K=35$ target classes, and $\mathbf{p} = \text{Softmax}(\mathbf{z})$.

\section{GPU Baseline Training Run and ONNX Export}
The baseline model was trained over 200 epochs on an **NVIDIA GeForce RTX 4050 Laptop GPU (6 GB VRAM, CUDA 12.1)** using the **AdamW optimizer** ($\text{lr} = 2 \times 10^{-3}$, $\text{weight\_decay} = 1 \times 10^{-4}$) and a **Cosine Annealing Learning Rate Scheduler**.

\begin{table}[htbp]
\centering
\caption{Baseline GPU Training Benchmark Metrics}
\label{tab:gpu_benchmarks}
\begin{tabular}{@{}ll@{}}
\toprule
\textbf{Metric / Parameter} & \textbf{Value} \\
\midrule
Dataset Size & 107,517 Clean 3D Hand Vectors \\
Train / Val / Test Split & 80\% (86,013) / 10\% (10,752) / 10\% (10,752) \\
Training Accuracy & \textbf{99.89\%} \\
Validation Accuracy & \textbf{99.67\%} \\
Held-Out Test Accuracy & \textbf{99.70\%} (10,720 / 10,752 Correct) \\
Training Time (200 Epochs) & 4 minutes 12 seconds \\
Exported ONNX Binary Size & \textbf{556 KB} (\texttt{isl\_letter\_classifier.onnx}) \\
ONNX CPU Inference Latency & \textbf{$<1.8\text{ ms}$ per frame} \\
\bottomrule
\end{tabular}
\end{table}

\section{The Inter-User Anatomical Variance Challenge}
While the baseline model achieved 99.70\% test accuracy, live testing revealed that individual signers possess subtle anatomical differences in finger length, knuckle flexibility, and resting hand posture. To ensure our system adapts to any individual user without requiring manual re-annotation, we engineered a personalization and geometric augmentation pipeline.

\section{Personalized Sign Recorder Studio GUI}
We built an interactive desktop GUI application (\texttt{prototype/sign\_recorder\_studio.py}) featuring a structured two-stage recording protocol:
\begin{itemize}
    \item \textbf{3-Second Red Preparation Window:} Allows the signer to position their hand in front of the camera lens.
    \item \textbf{3-Second Green Active Recording Window:} Automatically streams and records 90 continuous 126-dimensional landmark vectors at 30 FPS.
\end{itemize}
Using this studio, we captured personalized sign vectors for \textbf{33 classes} (letters A--Y and digits 1--9) saved in structured JSON arrays under \texttt{data/user\_recorded/}.

\section{3D Geometric Spatial Augmentation Mathematics}
To maximize generalization and prevent overfitting to a single camera position, our augmentation engine (\texttt{prototype/augment\_and\_train\_master\_letters.py}) applies synthetic 3D spatial perturbations:
\begin{enumerate}
    \item \textbf{3D Synthetic Euler Rotations ($\theta_x, \theta_y, \theta_z \sim \mathcal{U}(-18^\circ, +18^\circ)$):}
    \begin{equation}
        \mathbf{P}_{\text{aug}} = \mathbf{R}_z(\theta_z) \mathbf{R}_y(\theta_y) \mathbf{R}_x(\theta_x) \mathbf{P}_{\text{norm}}
    \end{equation}
    \item \textbf{Isotropic Scale Jitter ($s \sim \mathcal{U}(0.88, 1.12)$):} Simulates distance shifts relative to the lens.
    \item \textbf{Gaussian Joint Jitter ($\mathbf{\delta} \sim \mathcal{N}(0, \sigma^2)$ with $\sigma = 0.012$):} Simulates micro-tremors and camera sensor noise.
\end{enumerate}
This pipeline synthesized a master dataset of **246,104 augmented samples**.

\section{8-Second GPU Co-Training and Zero-Forgetting Fine-Tuner}
Our fine-tuning engine (\texttt{prototype/fine_tune_engine.py}) uses a **co-training mini-batch strategy** blending 20\% user-recorded augmented vectors with 80\% global baseline vectors. Executing on the RTX 4050 GPU, fine-tuning completes in **~8 seconds**, elevating live validation accuracy to **99.84\%--99.96\%** with zero catastrophic forgetting of the global sign alphabet.
